% Responsável: TODO(grupo) — escrita coletiva.
% Rubrica: peso 2, empatado com detecção. É a seção que mais vale por página.
% Enunciado, item 6.3: limitações da técnica, casos em que falhou, o que mudaria
% com mais tempo/dados.
\chapter{Discussão crítica}
\label{cap:discussao}

\section{O que funcionou}
\label{sec:funcionou}

TODO(grupo).

\section{O que não funcionou}
\label{sec:nao-funcionou}

TODO(grupo): esta seção não é opcional nem uma confissão de fracasso --- é o
principal item da rubrica. Casos concretos, não generalidades.

\section{Análise dos erros}
\label{sec:erros}

TODO(grupo): quais janelas foram classificadas errado, sob qual carga, com qual
severidade de defeito, e a hipótese para cada padrão de erro.

\section{Limitações}
\label{sec:limitacoes}

TODO(grupo): defeitos do CWRU são induzidos por eletroerosão, não por fadiga
natural; generalização para outra bancada; efeito da carga.

\section{O que faríamos com mais tempo ou mais dados}
\label{sec:trabalhos-futuros}

TODO(grupo).
